\section{Grupo objetivo definido por un corte conocido}\label{sec:corte}

Las secciones~\ref{sec:canal} y~\ref{sec:instrumentos} trataban de qué significan las covariables registradas. Esta sección y las dos siguientes toman las covariables como dadas y tratan de lo bien que se estiman los pesos y de cómo expresar la incertidumbre; esta empieza por la forma más habitual de definir el grupo objetivo.

En el caso típico, el grupo objetivo es un corte en una puntuación que la organización conoce, $S=D\cdot\ind\{s(X)\le c\}$. Entonces no hace falta un clasificador que distinga el grupo objetivo del canal donante: basta la propensión de canal ordinaria, con el corte impuesto a mano.

\begin{proposicion}[Pesos con corte conocido]\label{prop:corte}
Sea $S=D\cdot\ind\{s(X)\le c\}$, con $s$ medible y conocida, y sea $e$ una versión de $\PP(D=1\mid X)$. Bajo (A3), existe $k>0$ tal que
\[
\frac{dP}{dQ}(x)=k\,\ind\{s(x)\le c\}\,\frac{e(x)}{1-e(x)}\qquad Q\text{-casi seguro.}
\]
\end{proposicion}
\begin{proof}
Sea $\mu$ la ley de $X$ en la población. Como $s(X)$ es función de $X$, $\PP(S=1\mid X)=\ind\{s(X)\le c\}\,e(X)$. Por la definición de probabilidad condicionada, $P(A)=\int_A\ind\{s\le c\}\,e\dd\mu/\PP(S=1)$ y $Q(A)=\int_A(1-e)\dd\mu/\PP(D=0)$ para todo $A\in\A$, así que $P$ y $Q$ son absolutamente continuas respecto de $\mu$, con densidades proporcionales a $\ind\{s\le c\}\,e$ y a $1-e$. El conjunto $\{e=1\}$ es $Q$-nulo, y fuera de él, bajo (A3), la regla de la cadena da $dP/dQ=(dP/d\mu)/(dQ/d\mu)=k\,\ind\{s\le c\}\,e/(1-e)$, con $k=\PP(D=0)/\PP(S=1)$.
\end{proof}

La identidad es elemental, pero cambia el problema de estimación. El clasificador del Teorema~\ref{teo:odds} tiene que aprender un salto nítido en $\{s=c\}$, con los pocos casos del grupo objetivo como positivos. La propensión de canal suele ser una función suave y se estima con todas las unidades del canal premium. La constante $k$ se cancela al normalizar, y el diagnóstico de soporte del Teorema~\ref{teo:lebesgue} se lee en $\hat e$: falta soporte donde $\hat e\to1$ dentro de $\{s\le c\}$. La misma identidad, aplicada a las covariables ampliadas con una $U$ no observada, da el Lema~\ref{lem:calibracion} de calibración de la sección~\ref{sec:sensibilidad}.

\paragraph{Relación con la literatura.} La identidad no es nueva en sustancia. En el análisis por subgrupos con puntuaciones de propensión es práctica habitual estimar la propensión en toda la cohorte y aplicarla dentro del subgrupo \citep{rassen2012}, y \citet{dong2020} estudian cuándo conviene ese modelo y cuándo uno ajustado en cada subgrupo, que equilibra mejor dentro del subgrupo a costa de más varianza. Lo que añadimos es la comparación con el clasificador «grupo objetivo frente a donantes», el habitual cuando el cambio de medida se plantea como estimación de un cociente de densidades \citep{bickel2009,sugiyama2012}, en el caso en que el grupo objetivo es un corte en una puntuación de las propias covariables. Con modelos flexibles, los dos enfoques son consistentes y las diferencias son de muestra finita, pero afectan a cosas que la literatura de subgrupos no examina: el signo del recorte, el tamaño muestral efectivo y la validez del bootstrap (sección~\ref{sec:res-corte}).

\paragraph{Estimación.} En los estimadores de la sección~\ref{sec:estimadores}, $\hat o$ se sustituye por $\ind\{s\le c\}\,\hat e/(1-\hat e)$, y la muestra incluye también las unidades premium de fuera del grupo objetivo. Con muestreo poblacional de $(D,S,X,Y)$, la función de influencia correspondiente es $\PP(S=1)^{-1}\bigl[S(m_0-\theta)+(1-D)\,\ind\{s\le c\}\,e/(1-e)\,(Y-m_0)\bigr]$; esperamos un análogo del Teorema~\ref{teo:normalidad}, pero no lo demostramos.
